\providecommand{\Needspace}[1]{}
\Needspace{8\baselineskip}
\item \textbf{Limpieza y resumen de un catalogo estelar con pandas}.
\nopagebreak[4]

\begin{tcolorbox}[colback=blue!2!white,colframe=blue!55!black,title={Enunciado},
                  before skip=3pt,after skip=5pt,boxsep=3pt,top=3pt,bottom=3pt,
                  breakable]
\small
\raggedright
Un catalogo pequeno de estrellas fue armado desde distintas fuentes y guardado
en un archivo CSV. Antes de comparar las estrellas, se debe leer la tabla,
inspeccionarla, eliminar mediciones incompletas y crear una columna fisica
derivada.

El archivo CSV se encuentra en la siguiente ruta:

\begin{pythonjupyter}
url = (
    "https://raw.githubusercontent.com/jperaltac/pcfi161/"
    "main/Solemnes/2026/S3/catalogo_estrellas_s3.csv"
)
\end{pythonjupyter}

Realice las siguientes operaciones:

\begin{itemize}
\setlength{\topsep}{4pt}
\setlength{\itemsep}{3pt}
\setlength{\parsep}{0pt}
\item[(a)] Importe \texttt{numpy}, \texttt{pandas} y \texttt{matplotlib.pyplot}.
Lea el archivo CSV usando \texttt{pd.read\_csv(url)} y guardelo en un DataFrame
llamado \texttt{estrellas}.
\item[(b)] Inspeccione la tabla mostrando las primeras filas, dimensiones, tipos
de datos y cantidad de valores faltantes por columna.
\item[(c)] Cree un DataFrame llamado \texttt{estrellas\_limpias} eliminando solo
las filas que tengan datos faltantes en
\texttt{temperatura\_K}, \texttt{radio\_Rsol}, \texttt{distancia\_pc} o
\texttt{magnitud\_V}. Use \texttt{copy()} al final.
\item[(d)] En \texttt{estrellas\_limpias}, cree la columna
\texttt{luminosidad\_rel} usando
\[
L_{\rm rel}=R^2\left(\frac{T}{5778}\right)^4,
\]
donde \(R\) es \texttt{radio\_Rsol} y \(T\) es \texttt{temperatura\_K}.
\item[(e)] Seleccione las estrellas candidatas que cumplen simultaneamente:
\[
\texttt{distancia\_pc < 20},\qquad
\texttt{magnitud\_V < 0.5},\qquad
\texttt{luminosidad\_rel > 5}.
\]
Muestre solamente \texttt{nombre}, \texttt{tipo}, \texttt{distancia\_pc},
\texttt{magnitud\_V} y \texttt{luminosidad\_rel}.
\item[(f)] Agrupe \texttt{estrellas\_limpias} por \texttt{tipo}. Para cada tipo,
calcule cantidad de estrellas, temperatura media, distancia mediana y
luminosidad relativa media. Ordene el resumen de mayor a menor luminosidad
relativa media.
\item[(g)] Grafique \texttt{temperatura\_K} versus \texttt{luminosidad\_rel}
usando un color distinto para cada \texttt{tipo}. Agregue etiquetas, leyenda y
grilla.
\item[(h)] Imprima una frase que indique cual es el tipo con mayor luminosidad
relativa media. La respuesta debe salir desde la tabla agrupada, no escrita a
mano.
\end{itemize}
\end{tcolorbox}

\begin{solution}
\begin{pythonjupyter}
import pandas as pd
import numpy as np
import matplotlib.pyplot as plt

url = (
    "https://raw.githubusercontent.com/jperaltac/pcfi161/"
    "main/Solemnes/2026/S3/catalogo_estrellas_s3.csv"
)

estrellas = pd.read_csv(url)

print(estrellas.head())
print("Dimensiones:", estrellas.shape)
print(estrellas.dtypes)
print("Valores faltantes:")
print(estrellas.isna().sum())

columnas_necesarias = [
    "temperatura_K",
    "radio_Rsol",
    "distancia_pc",
    "magnitud_V"
]

estrellas_limpias = estrellas.dropna(
    subset=columnas_necesarias
).copy()

estrellas_limpias["luminosidad_rel"] = (
    estrellas_limpias["radio_Rsol"]**2 *
    (estrellas_limpias["temperatura_K"] / 5778)**4
)

mask = (
    (estrellas_limpias["distancia_pc"] < 20) &
    (estrellas_limpias["magnitud_V"] < 0.5) &
    (estrellas_limpias["luminosidad_rel"] > 5)
)

candidatas = estrellas_limpias.loc[
    mask,
    ["nombre", "tipo", "distancia_pc", "magnitud_V",
     "luminosidad_rel"]
]
print(candidatas)

resumen = estrellas_limpias.groupby("tipo").agg({
    "nombre": "count",
    "temperatura_K": "mean",
    "distancia_pc": "median",
    "luminosidad_rel": "mean"
})

resumen = resumen.rename(columns={"nombre": "cantidad"})
resumen = resumen.sort_values(
    "luminosidad_rel",
    ascending=False
)

print(resumen)

plt.figure(figsize=(7, 5))
for tipo in estrellas_limpias["tipo"].unique():
    sub = estrellas_limpias[estrellas_limpias["tipo"] == tipo]
    plt.scatter(
        sub["temperatura_K"],
        sub["luminosidad_rel"],
        label=f"tipo {tipo}"
    )

plt.xlabel("Temperatura [K]")
plt.ylabel("Luminosidad relativa")
plt.title("Temperatura versus luminosidad relativa")
plt.legend()
plt.grid(True, alpha=0.3)
plt.show()

tipo_mayor = resumen.index[0]
lum_mayor = resumen.iloc[0]["luminosidad_rel"]

print(f"El tipo con mayor luminosidad relativa media es {tipo_mayor}")
print(f"con L_rel media = {lum_mayor:.2f}.")
\end{pythonjupyter}
\end{solution}
